%  LaTeX support: latex@mdpi.com 
%  For support, please attach all files needed for compiling as well as the log file, and specify your operating system, LaTeX version, and LaTeX editor.

%=================================================================
\documentclass[data,article,submit,pdftex,moreauthors]{Definitions/mdpi} 

%=================================================================
%----------
% journal
%----------

% MDPI internal commands
\firstpage{1} 
\makeatletter 
\setcounter{page}{\@firstpage} 
\makeatother
\pubvolume{1}
\issuenum{1}
\articlenumber{0}
\pubyear{2026}
\copyrightyear{2026}
%\externaleditor{Academic Editor: Firstname Lastname}
\datereceived{07 May 2026} 
%\daterevised{} %
\dateaccepted{} 
\datepublished{} 
\hreflink{https://doi.org/} % If needed use \linebreak
\usepackage{tabularx}
\usepackage{booktabs}
\usepackage[flushleft]{threeparttable}
\usepackage{ltablex} % Combines tabularx and longtable
\usepackage{booktabs} % For \toprule, \midrule, etc.
\usepackage{array} % For custom column definitions
\usepackage{caption} % For better caption control

%=================================================================
% Add packages and commands here. The following packages are loaded in our class file: fontenc, inputenc, calc, indentfirst, fancyhdr, graphicx, epstopdf, lastpage, ifthen, lineno, float, amsmath, setspace, enumitem, mathpazo, booktabs, titlesec, etoolbox, tabto, xcolor, soul, multirow, microtype, tikz, totcount, changepage, attrib, upgreek, cleveref, amsthm, hyphenat, natbib, hyperref, footmisc, url, geometry, newfloat, caption

\graphicspath{{figures/}} % path to folder with figures
\usepackage{subcaption}
\usepackage{listings}
\usepackage{xcolor}
\usepackage{gensymb} % for \textdegree
\usepackage{tabularx}
\newcolumntype{Y}{>{\centering\arraybackslash}X}
\usepackage{amsmath,mathtools,amsthm}
	\setcounter{MaxMatrixCols}{15}
\usepackage{array}
\usepackage{amssymb}
\usepackage{amsfonts}
\usepackage{float}
\usepackage{chngcntr}
\counterwithin*{equation}{section}
\usepackage[super]{nth}
\usepackage{longtable}
\usepackage{tabularx}
\usepackage{multirow}
\usepackage{adjustbox}

\definecolor{deepblue}{rgb}{0,0,0.5}
\definecolor{deepred}{rgb}{0.6,0,0}
\definecolor{deepmagenta}{RGB}{195,12,214}
\definecolor{deepgreen}{RGB}{29,122,30}
\definecolor{mintgreen}{RGB}{192,249,192}
\definecolor{codepurple}{RGB}{204, 0, 153}
\definecolor{LightCyan}{rgb}{0.88,1,1}
\definecolor{GainsboroPurple}{RGB}{247,176,247}
\definecolor{LightSlateGrey}{RGB}{124,118,118}
%    
\lstdefinestyle{mystyle}{
    commentstyle=\color{LightSlateGrey},
    keywordstyle=\color{deepgreen}\bfseries,
    emphstyle=\color{deepred},
    numberstyle=\tiny\color{deepblue},
    stringstyle=\color{codepurple},
    basicstyle=\ttfamily\scriptsize,
    breakatwhitespace=false,         
    breaklines=true,                 
    captionpos=b,                    
    keepspaces=true,                 
    numbers=left,                    
    numbersep=5pt,                  
    showspaces=false,                
    showstringspaces=false,
    showtabs=false,                  
    tabsize=2
}
\lstset{style=mystyle}


%=================================================================
% Full title of the paper (Capitalized)
\Title{Artificial Intelligence and Big Data Analytics for Seismic Hazard Assessment: Advancements in the Marmara Region, Türkiye}

% MDPI internal command: Title for citation in the left column
\TitleCitation{Artificial Intelligence and Big Data Analytics for Seismic Hazard Assessment: Advancements in the Marmara Region, Türkiye}

% Author Orchid ID: enter ID or remove command
\newcommand{\orcidauthorA}{0000-0002-5759-1089} % Polina Add \orcidA{} behind the author's name

% Authors, for the paper (add full first names)
%\Author{Firstname Lastname $^{1}$\orcidA{}, Firstname Lastname $^{2}$ and Firstname Lastname $^{2,}$*}
\Author{Polina Lemenkova $^{1}$*\orcidA{}, Abdullah Can Zülfikar $^{1}$}

%\longauthorlist{yes}

% MDPI internal command: Authors, for metadata in PDF
\AuthorNames{Polina Lemenkova, Firstname Lastname}
% MDPI internal command: Authors, for metadata in PDF

% MDPI internal command: Authors, for citation in the left column
\AuthorCitation{Lemenkova, P., Zülfikar, A. C.} 

% MDPI internal command: Authors, for citation in the left column

% If this is a Chicago style journal: Lastname, Firstname, Firstname Lastname, and Firstname Lastname.

% Affiliations / Addresses (Add [1] after \address if there is only one affiliation.)
\address{%
$^{1}$ \quad Disaster and Emergency Management Department, Institute of Earthquake Engineering and Disaster Management, Istanbul Technical University. İTÜ Ayazağa, 34469, Maslak Kampüsü, İstanbul, Türkiye.
}

% Contact information of the corresponding author
\corres{Correspondence: lemenkova@itu.edu.tr}

% Abstract (Do not insert blank lines, i.e. \\) 
\abstract{The Eastern Marmara region of Türkiye represents a critical junction of the Eurasian, African, and Arabian plates, characterized by high seismic vulnerability along the North Anatolian Fault. Given the rapid urbanization of İstanbul and Kocaeli, the integration of advanced big data analytics is essential for effective hazard mitigation. This paper reviews recent advancements in geological and geomorphic risk assessment using large datasets, specifically focusing on the intersection of geophysical monitoring and Data Science. We evaluate the transition from traditional Geographic Information Systems (GIS) to AI-driven methodologies, driven by the explosion of seismic big data and remote sensing (RS) capabilities. The study systematically analyzes the application of Machine Learning (ML) and Deep Learning (DL) for geospatial data processing, including automated feature extraction in earthquake engineering and pattern recognition for landslide detection in satellite imagery. Core big data analytics techniques are reviewed, including the use of Random Forest, Support Vector Machines (SVM), and Gradient Boosting (XGBoost) for predictive modeling, alongside Convolutional Neural Networks (CNNs) for image segmentation and vulnerability mapping. Furthermore, we highlight the role of Python and R library ecosystems in facilitating statistical modeling and geophysical simulations. The synthesis of Natural Language Processing (NLP) for geospatial querying and AI-assisted cartographic visualization demonstrates a significant shift toward high-dimensional data integration in Earth observations. The findings suggest that AI-assisted data modeling significantly enhances the precision of seismic prognosis and infrastructure risk assessment in vulnerable mountain environments. By critically analyzing case studies within Türkiye’s tectonic context, this review demonstrates how advanced big data analytics transform raw geophysical observations into actionable preparedness strategies. This contribution underscores the vital role of interdisciplinary data management in improving the resilience of seismically active urban zones.}

% Keywords
\keyword{big data; seismology; cloud computing; geophysics; AI; data analysis; signal processing} 

% The fields PACS, MSC, and JEL may be left empty or commented out if not applicable
%\PACS{J0101}
%\MSC{}
%\JEL{}

\begin{document}

%----------- section ----------->
\section{Introduction}

%----------- section ----------->
\subsection{Background}

Seismic hazard monitoring systems prioritize the early detection and forecasting of earthquake events to mitigate geotechnical risks \cite{Allen2012,Kanamori2005}. Modern geophysical centers analyze thousands of global sensor streams per second, generating millions of daily data points \cite{USGS2020,IRIS2018,Chung2020}. However, the transition to real-time seismic forecasting necessitates processing capabilities that exceed the limits of traditional geophysical analysis \cite{Beroza2018,Allen2019}. The integration of Artificial Intelligence (AI) and cloud computing has transformed this landscape, enabling the treatment of global seismic signals as cohesive big data massifs rather than isolated regional records \cite{Perol2018,Zhu2019,Kong2019}.

The paradigm shift from traditional geophysics to big data seismology is not merely a matter of increased data volume; it represents a fundamental change in analytical methodology \cite{Bergen2019,Donner2019}. Contemporary seismic datasets require novel frameworks for storage and high-throughput processing to extract actionable prognostic insights \cite{Ahern2018,IRISDMC2015}. Understanding big data in this context involves recognizing that the complexity and scale of modern sensor networks demand a departure from legacy data management systems toward distributed, scalable data science architectures, Figure \ref{fig01}.

\begin{figure}[H]
    \begin{center}
	\hspace{-35pt}\includegraphics[width=13.0cm]{Fig_01.png}
    \end{center}
    \caption{Big Data in Geophysics and Seismology: 5Vs. Diagram source: authors.}
    \label{fig01}
\end{figure}

The "Five Vs" of big data—volume, velocity, variety, veracity, and value—define the current advantages over regional geophysical databases \cite{Laney2001,Kitchin2014,Hashem2015}. While traditional archives typically manage gigabytes of localized information, modern big data systems routinely ingest petabytes of global sensor data \cite{Ahern2018}. Furthermore, the velocity of these systems enables the processing of real-time data streams, providing the immediate latency-critical insights required for seismic risk evaluation and rapid decision-making \cite{gandomi2015beyond, hashem2015rise, allen2009earthquake}.

The variety in big data in geophysics is characterised by a spectrum of heterogeneous data structures \cite{li2016big, wang2019geophysical}. The complexity of formats arises due to the diverse data structures and sources of data that should be integrated and analysed together \cite{chen2014big}. Even if the data is available as a CSV in clear rows and columns, that file often presents a wrapper for a mix of information pulled from across geophysical sensors \cite{ma2015big,kessler2019data}. Finally, the veracity of large volumes of data is related to data quality and integrity \cite{fang2015big,reyes2017data}. When dealing with web-sourced datasets on seismic records, these might include incomplete signals, duplicate entries of events, inconsistent formatting of diverse data, occasionally broken sensors, or outdated information \cite{mousavi2019cred}.

The five Vs create interconnected challenges in modern seismology that must be addressed. Conventional methods of seismic data analysis mostly rely on manual examination of datasets, ignoring massive volumes of historical seismic data available at archives and new records flowing through sensors and stored in online geophysical repositories \cite{kong2019machine}. The analysis of such large massifs of data helps discovering anomalies in mantle processes with better precision and speed. Such insight can be derived from analysing real-time data characterising location and magnitude of seismic events, unstructured cartographic data, and all relevant geological characteristics that define the unique value of regional data \cite{perol2018convolutional,zhu2019phasenet}. As a result, big data analytics in seismology prevent failures in prognosis through predictive analytics that processes sensor data from thousands of machines simultaneously.

%----------- section ----------->
\subsection{Research problem}

Big data revolutionise geophysics through a possibility of dynamic risk assessment. For instance, in Seismic Hazard Risk Assessment (SHRA), big data has shifted the field from static, once-a-decade maps to dynamic, high-resolution models of real-time monitoring. Data integration and real-time processing enables historical data analysis, current data monitoring and future data forecasting \cite{Sen2009}. Recent decades have witnessed diverse geological and climate hazards which underscore their cumulative effects on the environment (Fan et al., 2019; McGuire, 2013). The ongoing influence of geological activities, such as earthquakes, coupled with climate-related events like extreme weather patterns and rising sea levels, reflects a complex interplay that warrants further examination \cite{Liggins2013}. Data-driven analysis helps understanding the nature of these hazards to develop strategies of risk mitigation \cite{DeSalvio2023,Glasgow2021}. Seismic risks are particularly prominent due to their devastating consequences for society, necessitating effective seismic risk assessment methods \cite{Dowrick2009,Liu2024}. 

Rapid data analysis is essential in earthquake monitoring. While traditionally seismology relied on a few expensive, high-grade stations, nowadays, it is characterised by leverage an explosion of data from diverse sources to predict how the ground will shake and how the built environment will respond. This is especially important for Marmara Sea region of Türkiye, which is notably vulnerable to seismic activities, Figure \ref{fig02}. 

\begin{figure}[H]
    \begin{center}
	\hspace{-35pt}\includegraphics[width=13.0cm]{Fig_02.png}
    \end{center}
    \caption{Geophysical Data Processing Framework: Istanbul EEW \& Seismic Analysis. Diagram source: authors.}
    \label{fig02}
\end{figure}

The country is located on the active tectonic faults, which contribute to the initiation of multiple earthquakes \cite{Burton1984,Yao2024}. Given this context, big data monitoring is essential for seismic risk assessment of the country, underscoring the need for ongoing research and the development of methodologies that can enhance societal safety.

%----------- section ----------->
\subsection{Objectives and goals}

In this paper, we analyse the challenge of big data technologies in geophysics through understanding how data impacts modern decisions of seismic hazard risk assessment. Employing these fundamentals in seismology support meaningful and data-driven forecasting and making informed recommendations about possibility of earthquakes. We discuss the challenge of SHRA in the perspective of big data and high-scale geophysical analytics which can bridge the gap between traditional data analysis in seismology and predictive analytics and prognostic forecasting of earthquakes. Most traditional methods of seismic data analysis use management of a CSV. In contrast, big data can architect an SHRA system that handles volumes, velocity and variety of seismic signals simultaneously for predictive analytics in applied geophysics. 

%----------- section ----------->

\section{Seismic Vulnerability in Türkiye: From Tectonic Plate Interactions to Historical Disaster Impacts}

\subsection{Geology and Tectonics in Türkiye: General Overview}

Türkiye is notably vulnerable to seismic activities, especially in regions near active tectonic faults, which contribute to the initiation of earthquakes \cite{Basaran2014,Kandregula2024}. Given this context, geological monitoring is essential for comprehensive seismic risk assessment, underscoring the need for ongoing research and the development of methodologies that can enhance societal safety \cite{Korkmaz2013}. In Türkiye, the regions of seismically active areas are situated at the border of the Anatolian Block \cite{Jolivet2023}. These areas are located within a zone of lithosphere collision, which is characterized by the interaction of three major tectonic plates: the Eurasian Plate, the African Plate, and the Arabian Plate.

This tectonic activity significantly influences the geological and seismic characteristics of the region. The Anatolian Block is being squeezed westward as the Arabian Plate moves northward against the Eurasian Plate, while the African Plate subducts beneath the Hellenic Arc \cite{McKenzie1972, LePichon2010}. Such processes, which influence regional tectonics, contribute to the occurrence of frequent and powerful earthquakes at lithosphere boundaries. The most prominent feature of this system is the North Anatolian Fault Zone (NAFZ), a 1,200 km right-lateral strike-slip fault extending from the Gulf of Saros to Karlıova \cite{Barka1992, Sengor2005}. It was formed 13–11 million years ago and remains seismically active, capable of producing significant earthquakes up to magnitude 8.0 \cite{Stein1997}. Current deformation, particularly in the Marmara region, is asymmetric and heavily influenced by the complex regional geology reflecting conditions along the NAFZ \cite{Puccinelli2021}. This ongoing deformation ensures that the region remains one of the most monitored and researched seismic zones in the world, as the accumulation of strain along these boundaries poses a continuous threat to urban centers.

\subsection{Historical Outline of Seismic Events in Türkiye (496 BC to Present)}

The following historical overview examines the relentless cycle of seismic activity in the Anatolian region, tracing a path from the earliest documented tremors of antiquity through the devastating tectonic ruptures of the modern era, Figure \ref{fig03}. 

\begin{figure}[H]
    \begin{center}
	\hspace{-35pt}\includegraphics[width=14.0cm]{Fig_03.png}
    \end{center}
    \caption{Historical outline of known major seismic events in Türkiye. Diagram source: authors.}
    \label{fig03}
\end{figure}

The historical record of the Anatolian region reveals a relentless cycle of seismic activity that has shaped human settlement for millennia. In the Aegean region, specifically the Gulf of Izmir, the Karaburun peninsula, and the island of Chios, historical records indicate that from 496 BC to 1949 AD, the area experienced a total of 20 moderate to severe earthquakes \cite{altinok2011, papazachos2000}. Many of these ancient events resulted in devastating tsunamis, with notable occurrences recorded in the years 1389, 1856, 1866, 1881, and 1949 \cite{yalciner2002}.

During the 19th century, the region faced intensified activity. The 1881 Chios-Çeşme earthquake stands out as a catastrophic event that caused thousands of fatalities and significant coastal changes \cite{ambraseys2000}. Moving into the 20th century, the North Anatolian Fault Zone began a well-documented westward-migrating sequence of large ruptures. A significant event in this timeline was the Amorgos earthquake in 1956, which struck the Southern Aegean \cite{okals2009}. This earthquake generated a tsunami with a wave height of 2.5 meters, leading to severe flooding over an area with a depth of 1.5 kilometers on the island of Kalimnos. Additionally, this seismic activity caused substantial material damage to several nearby islands, including Ios, Nisyros, Patnos, Tilos, and Lipsos \cite{papadopoulos2007}.

The tsunami wave reached 1m in Fethiye, and the flood depth was observed to be 250m.On 17 August 1999, the $M_w$ 7.4 İzmit earthquake occurred, as documented by Inan et al. in 2007. This seismic event was significant in terms of its magnitude and impact, reflecting the geological vulnerabilities of the region. Notably, the 1999 earthquakes heightened stress in the Marmara segment, where a 50\% chance of a magnitude 7.6 event is anticipated between 2005 and 2055. On July 21, 2017, an earthquake with a magnitude of $M_w$=6.6 occurred in the Gulf of Gökova at a depth of approximately 6 km, triggering a tsunami that reached a maximum run-up height of 1.9 m in Bodrum.

Most recently, on February 6, 2023, Türkiye experienced a catastrophic seismic sequence characterized by a mainshock of 7.8 $M_w$, followed by two significant earthquakes measuring 7.7 and 7.6 $M_w$. The epicenters were located in Pazarcık and Elbistan within the Kahramanmaraş province, with another strike-slip shock located 37 kilometers northwest of Gaziantep. This event reached a Mercalli intensity of XII (Extreme), causing total destruction in Antakya and surrounding provinces. As documented by \cite{Akcan2024} and \cite{Akinci2025}, these events represent some of the most powerful shocks in the region’s modern history.

\subsection{Seismic Risk Assessment and the Need for EEW}

The extensive destruction and widespread damage observed across Adıyaman and adjacent provinces in 2023 serve as a stark reminder of the far-reaching consequences of these hazards \cite{Hussain2023,Yilmaz2025,Toprak2025,Carena2023}. Such devastating geologic and seismic hazards have profound impacts on both urban infrastructure safety and the natural landscape. The damages inflicted on urban areas can significantly hinder social-economic development and adversely affect human well-being. This interconnectedness highlights the necessity for robust urban planning and infrastructure resilience strategies and initiatives in Türkiye. 

Specifically, there is an urgent need for the development of Earthquake Early Warning (EEW) methodologies and methodologies that can enhance societal safety. Ongoing research into seismic risk assessment is essential to analyze the impact of earthquakes and improve future preparedness. By integrating real-time geological monitoring with rapid response strategies, the risk to communities can be mitigated. The goal of current and future research is to transform our understanding of the Anatolian Block’s tectonic behavior into actionable safety protocols that protect urban centers from the inevitable recurrence of major seismic events.


\section{Research gap identification and contribution}

To the best of our knowledge, no comprehensive review papers exist with a focus on developing methods of big data analysis for rapid evaluating of seismic risks in Türkiye. Despite the reported individual case studies, the integration review analyses comprising big data issues and seismic monitoring are still missing due to high complexity of this inter-disciplinary domain. To fill in this gap, this review article provides a comprehensive analysis of the evolution of techniques for assessing big data perspectives for assessment of geologic and seismic risks in Türkiye. It conducts an in-depth evaluation of existing case studies to highlight advancements in AI methodologies of big data processing and their practical applications in the context of regional risk assessment of Türkiye. Mitigating earthquake risks can be achieved by implementing an EEW system alongside consistent seismic monitoring using big data, Figure \ref{fig04}. Monitoring and preventive measures of EEW are crucial for enhancing public awareness and minimizing potential damages.

\begin{figure}[H]
    \begin{center}
	\hspace{-35pt}\includegraphics[width=14.0cm]{Fig_04.png}
    \end{center}
    \caption{Development of Big Data in Seismology, EEW \& Geophysical Monitoring in Istanbul, Türkiye. Diagram source: authors.}
    \label{fig04}
\end{figure}

\subsection{Current needs and challenges}
Today, geophysics operates at a complex intersection where the volume of data is increasing, but the physics-based models used to interpret that data are hitting computational and theoretical ceilings \cite{Gueneyli2026}. As industry focuses shift toward high-density data acquisitions for seismic signals or geothermal exploration, the scale of the data has made traditional manual workflows impossible. Hence, modern geophysics faces a data deluge from high-density nodal sensors that overwhelms traditional Input-Output (I/O) bandwidth and storage, making real-time "edge" processing essential but difficult to implement \cite{Bacon2025,Seydoux2024}.

While powerful, the legacy toolkit of geophysics and existing approaches of seismic data processing has several inherent bottlenecks and limitations, Table \ref{tab01}.

\begin{table}[ht]
\centering
\begin{threeparttable}
\caption{Limitations of Existing Geophysical Processing Approaches}
\label{tab01}
\begin{tabularx}{\textwidth}{@{} l l X @{}}
\toprule
\textbf{Approach} & \textbf{Core Constraint} & \textbf{Big Data \& Geophysical Impact} \\ \midrule
Traditional Inversion & Non-uniqueness & Multiple subsurface models satisfy the same observed potential field or seismic data (Equivalence problem). High-dimensionality in big data exacerbates uncertainty quantification. \\ \addlinespace
Linearized Models & Oversimplification & Standard kernels assume a linear Earth. Real complex geology (e.g., salt tectonics) is highly non-linear; forcing linearity in high-volume datasets results in "false anomalies" and signal artifacts. \\ \addlinespace
Manual QC & Human Bottleneck & High-density 3D/4D seismic surveys generate petabytes of traces. It is physically impossible for human interpreters to review every trace, making manual QC the primary latency in the processing pipeline. \\ \addlinespace
HPC Constraints & Data Movement & In seismic migration (e.g., RTM), moving massive velocity models from storage to compute nodes consumes more energy than the FLOPs performed (Von Neumann bottleneck\tnote{1}). \\ \bottomrule
\end{tabularx}
\begin{tablenotes}
\small
\item[1] The Von Neumann bottleneck refers to the limited throughput between the central processing unit (CPU) and memory compared to the amount of memory.
\end{tablenotes}
\end{threeparttable}
\end{table}

The primary technical hurdle lies in integrating physics-based wave equations with AI. While deep learning (DL) offers computational efficiency, it often lacks the geological rigor of physics-driven inversion, which remains inherently ill-posed and affected by non-uniqueness, where multiple subsurface models can explain the same data \cite{Li2019}. Furthermore, conventional approaches are limited by the trade-off between resolution and depth of investigation and by strong nonlinearity in complex settings such as sub-salt environments, while their reliance on iterative workflows and manual quality control introduces scalability bottlenecks for modern large-scale datasets \cite{Brossier2015}.

%----------- section ----------->
\section{Big data processing in geophysics}

%----------- subsection ----------->
\subsection{Hybrid computing: HPC and Cloud mosaic}
One of the most demanding data processing tasks in geophysics includes those involving the massive datasets from nodal sensors and the complex inversions, such as Full Waveform Inversion (FWI) \cite{witte2019scalable}. Since nodal datasets can reach petabyte scales, geophysicists no longer move all data to a single location; instead, they use a hybrid mosaic approach \cite{peter2011forward, wang2020cloud}. The principle of hybrid approach consists in taking the advantages of both HPC and cloud computing \cite{fox2017big, assuncao2015big}.

Handling the big data in geophysics requires a shift from linear, sequential processing to a hybrid architecture that blends high-performance computing (HPC) with cloud-native intelligence \cite{fox2017big, hashem2015rise}. To process massive nodal datasets and complex inversions, geophysicists utilise a hybrid architecture that combines the raw power of GPU-accelerated high-performance computing (HPC) clusters with the elastic scalability of the cloud, effectively moving the processing code to the data to avoid massive transfer bottlenecks \cite{wang2020cloud}. For example, FWI is a compute-intensive task which simulates seismic waves through every cubic meter of the earth and requires heavy lifting on HPC. The hybrid infrastructure accelerates high-resolution FWI by embedding physical laws directly into AI algorithms, since it is driven by Physics-Informed Neural Networks (PINNs) \cite{raissi2019physics, sun2020surrogate}.

The FWI remains on GPU-accelerated HPC clusters because of the tight coupling required between processors. In turn, bursty tasks, such as initial seismic signal denoising and geometry correction of the dataset are offloaded to the cloud which profits from its elasticity. This allows teams to spin up 10,000 nodes for a single weekend of processing without owning the hardware. For remote nodal arrays, smart nodes now perform initial data compression and basic quality checks at the source, transmitting only the most relevant seismic signal back to headquarters, Table \ref{tab02}. 

% Set the tabularx environment to behave like a longtable
\keepXColumns 

\begin{tabularx}{\textwidth}{@{} >{\hsize=0.8\hsize\raggedright\arraybackslash}X >{\hsize=1.2\hsize\raggedright\arraybackslash}X >{\hsize=1.0\hsize\raggedright\arraybackslash}X @{}}
    \caption{State-of-the-Art Big Data Processing in Geophysics (2026 Framework)} \label{tab02} \\
    \toprule
    \textbf{Geophysical Task} & \textbf{2026 Advanced Strategy} & \textbf{Technology Stack} \\ 
    \midrule
    \endfirsthead   
    % Header for subsequent pages
    \multicolumn{3}{c}{{\bfseries \tablename\ \thetable{} -- Continued from previous page}} \\
    \toprule
    \textbf{Geophysical Task} & \textbf{2026 Advanced Strategy} & \textbf{Technology Stack} \\ 
    \midrule
    \endhead    
    % Footer for the bottom of the first page
    \midrule
    \multicolumn{3}{r}{{Continued on next page...}} \\
    \endfoot    
    % Final footer
    \bottomrule
    \endlastfoot
    % Table Data
    \textbf{High-Density Nodal Ingest} & Decoupled asynchronous streaming from 100k+ channel nodes; real-time metadata indexing. & Azure Blob/S3, Apache Kafka, Zarr/ASDF Formats \\ \addlinespace
    
    \textbf{Elastic FWI \& Velocity Modeling} & Multi-parameter PINN-based inversion to resolve complex salt geometries and anisotropy. & NVIDIA H200 Clusters, PyTorch/JAX-Seismic \\ \addlinespace
    
    \textbf{Seismic Signal Enhancement} & Self-supervised Deep Learning Autoencoders for 5D interpolation and ghost reflection removal. & Transformer-based Denoising, GANs \\ \addlinespace
    
    \textbf{Microseismic Monitoring} & Automated event detection and location using Edge-AI for real-time hydraulic fracture mapping. & NVIDIA Jetson AGX Orin, 5G/6S Telemetry \\ \addlinespace
    
    \textbf{4D Reservoir Characterization} & Digital Twin synchronization integrating seismic, EM, and production data for predictive flow modeling. & Azure Digital Twins, NVIDIA Omniverse \\ \addlinespace
    
    \textbf{Multi-Physics Integration} & Joint inversion of gravity, magnetic, and seismic data via Large Language Models (LLMs) for interpretation. & Foundation Models, Graph Neural Networks (GNNs) \\ 
\end{tabularx}


As a result of hybrid approach, geophysicists integrate edge computing for real-time data compression and automated AI quality control for trace picking. Such approach ensures geological accuracy while bypassing the "non-uniqueness" of traditional methods. Hence, the industry has shifted from labor-intensive manual workflows to an asynchronous, "digital twin" approach that can reconstruct complex subsurface structures with unprecedented speed and precision \cite{}.

%----------- subsection ----------->
\subsection{Architecture Paradigms for Marmara Regional Seismology}

To manage the heterogeneous data demands of the Marmara region, cluster architectures are categorized into three patterns optimized for distinct geophysical workflows. The Master-Worker paradigm, utilizing Apache Spark, facilitates high-precision batch analytics and task coordination essential for computationally intensive Full Waveform Inversion (FWI) and 3D tomographic modeling of the North Anatolian Fault. For high-density nodal arrays, a Peer-to-Peer pattern (e.g., Apache Cassandra) employs decentralized gossip protocols and N-replication to ensure maximum fault tolerance and high availability for seismic traces, eliminating single points of failure during regional events. Finally, a Hybrid Broker-Based architecture using Apache Kafka provides the decoupled scaling necessary to ingest and partition massive telemetry streams from borehole and GPS sensors. This framework balances the sub-second latency required for EEW against the high-velocity throughput needed for petabyte-scale tectonic research, Figure \ref{fig04}

%----------- subsection ----------->
\subsection{Physics-Informed Neural Networks (PINNs)}

The most demanding inversions are now handled by PINNs \cite{Aghamiry2021}. Unlike traditional AI, which might suggest geologically impossible structures, PINNs embed the actual wave equations (the physics) into the neural network's loss function. The Benefit of using PINNS is that they can "fill in the gaps" of sparse nodal data by ensuring every predicted pixel adheres to the laws of thermodynamics and fluid dynamics \cite{RashtBehesht2022,Moseley2020}. Moreover, AI models like PhaseNet (e.g., see \cite{ZhuBeroza2019,Park2020,Ni2023}) are used to automatically "pick" millions of seismic arrivals from nodal data—a task that would take human teams years to complete manually.

%----------- subsection ----------->
\subsection{Asynchronous Workflows and Compressed Sensing}

By 2026, the industry has largely solved the "bottleneck" problem by moving the code to the data (Cloud) rather than the data to the code, and by using AI as a "physics-aware" accelerator rather than just a pattern matcher. For instance, to manage the velocity of big data from fiber-optic Distributed Acoustic Sensing (DAS) and massive nodal arrays, current geophysics suggests the use of two methods:
\begin{itemize}
	\item Compressed sensing enables the reconstruction of high-resolution images from significantly fewer measurements by exploiting the inherent sparsity of signals or geological structures \cite{Candes2006,Donoho2006,Herrmann2012}.
	\item Asynchronous inversion approaches process large datasets in sequential or sliding temporal windows rather than waiting for complete survey acquisition, enabling near-real-time updates of subsurface models during active monitoring \cite{Gineste2021,Wilkinson2022,Minson2014}.
\end{itemize}

%----------- subsection ----------->
\section{Cluster computing suitability assessment}

%----------- subsection ----------->
\subsection{Distributed data processing}

Distributed data processing is most critical for geophysical projects. It is characterized by high-density data volumes and the need for iterative, high-fidelity physics simulations. High-resolution Carbon Capture and Storage (CCS) monitoring benefits immensely, as the 4D seismic workflows required to detect subtle $CO_2$ leaks demand parallelized comparison of petabyte-scale historical and real-time datasets. Similarly, Geothermal Site Characterization relies on GPU-accelerated clusters to perform Full Waveform Inversion (FWI) through complex, fractured basement rock; this process requires thousands of wave-propagation iterations that are mathematically impossible to solve on localized hardware within a viable timeframe. 

Projects involving Distributed Acoustic Sensing (DAS) also require these architectures to manage the extreme data velocity generated by fiber-optic cables, using edge clusters for real-time denoising and event detection. Finally, creating regional-scale Crustal Digital Twins for earthquake early-warning systems necessitates a distributed approach to run unsupervised AI models across thousands of nodes simultaneously. These scenarios illustrate that distributed computing is no longer a luxury but a necessity for projects where the spatial resolution of the survey and the physical complexity of the earth model exceed the memory and processing limits of monolithic systems, enabling faster risk mitigation and more precise subsurface imaging.

%----------- subsection ----------->
\subsection{Distributed computing frameworks}

Distributed computing frameworks are designed to process massive datasets such as big data in high speed. One example of such toolkits in open-source used in geophysics is PySpark -- the Python API for Apache Spark, Figure \ref{fig05}. 

\begin{figure}[H]
\centering
\includegraphics[width=13.0cm]{Fig_05.png}
 \caption{Integrated EEW \& Regional Seismic Data Infrastructure in Distributed computing frameworks in Istanbul. Diagram source: authors.}
 \label{fig05}
\end{figure}

In the context of earthquake monitoring and seismic signal processing, distributed computing frameworks such as PySpark enable the transition from analyzing isolated events to processing continuous waveform data across large-scale sensor networks \cite{Armbrust2018,Zaharia2016}. By leveraging DataFrame-based transformations to window high-frequency time series and querying historical seismic catalogs, geophysicists can identify subtle precursory patterns—such as tectonic tremor and seismic swarms—that are difficult to detect with traditional single-node systems \cite{Obara2002,Ide2007}. This distributed approach is essential for collaborative seismic monitoring, enabling the synchronization of real-time data streams with advanced signal-to-noise ratio enhancement techniques.

Furthermore, scalable data-processing frameworks support the deployment of deep learning models such as recurrent neural networks (RNNs) for seismic phase picking and time-series analysis \cite{Mousavi2019,ZhuBeroza2019}, as well as the large-scale cross-correlation computations required for ambient noise interferometry \cite{Bensen2007}.

The primary collaboration challenges stem from the “Domain vs. Code” gap, where specialized physical algorithms must be translated into scalable implementations, and the technical friction of handling legacy data formats such as SEG-Y alongside modern cloud-native formats \cite{Hall2019,OpenVDS2019}. These teams often face “black-box” skepticism, where scientific interpretability conflicts with opaque machine learning models \cite{Rudin2019}, as well as resource contention due to the significant computational demands of large-scale 3D seismic inversion \cite{Virieux2009}. Using a unified platform such as Azure Databricks helps mitigate these issues by enabling collaborative, scalable data engineering and distributed computing workflows \cite{Armbrust2018a,Databricks2020}.

\subsection{Geophysical data hubs in Türkiye: Distributed architecture, collaboration and technical silo}
The distributed approach to data processing in earthquake monitoring centres corresponds to its current organisational structure of geophysical data centres. Since Türkiye is home to some of the most advanced geophysical data centres in the Mediterranean and Middle East, its primary focus is on seismic monitoring due to the country’s unique position on major fault lines. Data centres include a network of centres processing and analysing seismic and earthquake data. The following institutions as the primary hubs for geophysical data collection, processing, and distribution, Table \ref{tab03}:

% Required in preamble: \usepackage{ltablex, booktabs, caption}

\keepXColumns
\begin{tabularx}{\textwidth}{@{} 
    >{\hsize=0.6\hsize\raggedright\arraybackslash}X 
    >{\hsize=0.8\hsize\raggedright\arraybackslash}X 
    >{\hsize=1.6\hsize\raggedright\arraybackslash}X 
@{}}
    \caption{Principal Organizations and Big Data Activities in Marmara Region Seismology and EEW} \label{tab03} \\
    
    \toprule
    \textbf{Data Centre} & \textbf{Primary Focus} & \textbf{Seismic Big Data \& EEW Activities} \\ \midrule
    \endfirsthead
    
    % Header for subsequent pages
    \multicolumn{3}{c}{{\bfseries \tablename\ \thetable{} -- Continued from previous page}} \\
    \toprule
    \textbf{Data Centre} & \textbf{Primary Focus} & \textbf{Seismic Big Data \& EEW Activities} \\ \midrule
    \endhead
    
    % Footer for the bottom of pages before the last one
    \midrule
    \multicolumn{3}{r}{{Continued on next page...}} \\
    \endfoot
    
    % Final footer at the very end of the table
    \bottomrule
    \endlastfoot

    % Table Body
    \textbf{KOERI (Kandilli)} \newline \textit{Est. 1868 (Imperial Observatory)} & Operates the National Earthquake Monitoring Center (NEMC). & Processes real-time streams for the Istanbul EEW system, utilizing borehole and surface sensors to trigger automated gas/rail shutdowns in the Marmara region. \\ \addlinespace
    
    \textbf{AFAD} \newline \textit{Est. 2009 (Consolidating TEMAD/EIE)} & Manages the TDVMS (Earthquake Data Center System) with 1,100+ stations. & Acts as the official authority for high-velocity acceleration data used in real-time structural health monitoring and regional intensity mapping. \\ \addlinespace
    
    \textbf{MTA} \newline \textit{Est. 1935 (Ankara)} & Curates the Active Fault Map of Türkiye. & Big data activities focus on high-resolution GIS integration of paleoseismological data and surface rupture mapping to inform long-term probabilistic hazard models. \\ \addlinespace
    
    \textbf{TPAO} \newline \textit{Est. 1954} & Manages petabyte-scale 3D/4D seismic reflection catalogs. & Utilizes high-performance computing (HPC) for deep-water imaging in the Black Sea/Marmara, providing critical crustal structure insights via legacy SEG-Y archives. \\ \addlinespace
    
    \textbf{TÜBİTAK MAM} \newline \textit{Est. 1972 (Gebze)} & Focuses on marine geophysics and R\&D. & Processes high-resolution bathymetry and sub-bottom profiler data from the Marmara Sea to identify active submarine fault segments and liquefaction risks. \\ \addlinespace
    
    \textbf{Belbaşı Monitoring Center} \newline \textit{Est. 1951} & Specialized CTBT facility. & Analyzes global seismic waves for nuclear test detection; its high-precision sensors contribute to deep-crustal noise characterization and signal processing research. \\ \addlinespace
    
    \textbf{ITU Geophysics} \newline \textit{Est. 1952 (Department)} & Houses the N. Canıtez Lab. & Advanced academic processing of electromagnetic and seismic data; pioneers AI-driven feature extraction for Marmara-specific site-response and tomography. \\ \addlinespace
    
    \textbf{METU Earth Systems} \newline \textit{Est. 2002} & Integrates multi-disciplinary big data. & Primarily blending GNSS/GPS crustal deformation records with seismic catalogs to quantify the slip rate and strain accumulation along the North Anatolian Fault. \\ \addlinespace
    
    \textbf{Dokuz Eylül (DAUM)} \newline \textit{Est. 1987 (İzmir)} & Primary data hub for the Aegean region. & Conducts massive fault analysis and seismic cataloging to study complex graben systems and volcanic-seismic interactions. \\ \addlinespace
    
    \textbf{Kocaeli University} \newline \textit{Post-1999 expansion} & Center for urban geophysics. & Processes massive site-response and microzonation datasets for the Marmara industrial corridor, focusing on soil-structure interaction during high-amplitude shaking. \\ 

\end{tabularx}

A geophysical data centre is organised into a structure consisting of data engineers managing high-volume pipelines, geophysicists that provide domain-specific interpretation, and data scientists applying predictive modeling. The collaboration within the tri-layer structure of a geophysical data centre is primarily hindered by technical and semantic silos that emerge when sharing massive, complex datasets. Friction arises from data format incompatibility, where data engineers prioritize high-performance cloud formats like Parquet, while geophysicists remain tethered to legacy SEG-Y files, leading to costly "data wrangling" cycles. Additionally, code reproducibility is challenged by the gap between exploratory research notebooks and standardized production pipelines, often exacerbated by a lack of shared version control for "ground truth" datasets. These hurdles are compounded by a communication barrier regarding model interpretability, where the data scientist's focus on statistical accuracy may clash with the geophysicist’s requirement for physical consistency and scientific rigor.

Current sharing practices in geophysical data centres often rely on a fragmented mix of physical data duplication, static reporting, and disconnected code repositories, which inadvertently reinforces technical silos. Data is frequently shared through manual "copy-paste" exports from legacy desktop software into shared network drives, leading to version chaos and the loss of critical metadata. 

Code sharing often lacks centralized version control of scripts developed for seismic signal processing. For instance, specialized scripts may reside in isolated repositories or being exchanged via email, while main results are typically communicated through static PDFs that decouple the final insights from the underlying Spark logic. These silos are further pronounced by the use of "closed" ecosystem tools and manual ETL processes that trap data in project-specific formats, though the industry is increasingly shifting toward "Lakehouse" architectures and unified governance to enable real-time, zero-copy collaboration

\section{Databricks: Cloud-native geophysics}

\subsection{Scaling seismic resilience: insights in Türkiye}
Databricks as a cloud-based data engineering and analytics platform for processing massive datasets is actively used in geophysics. In particular, it is employed for large-scale seismic data processing, earthquake monitoring, and oil and gas exploration. For instance, its use is driven by the industry’s shift toward "Cloud-Native Geophysics," where the traditional "moving files" approach is being replaced by distributed computing on a Data Lakehouse. Seismologists use Databricks to build automated ETL pipelines for real-time seismic networks.

Türkiye is a primary hub for the big data architecture in geophysics due to its high seismic activity and the massive scale of data. The large data massifs are generated by data centres such as AFAD (Disaster and Emergency Management Authority) and KOERI (Kandilli Observatory and Earthquake Research Institute). To ensure both operational stability and scientific agility in geophysical data centres, the Databricks workspace adopt a hybrid structure that separates concerns by environment (Dev/Staging/Prod) at the top level while organizing internal content by project. For example, by using a Unified Data Lakehouse, a researcher in Ankara can run a PySpark script on the same data that a technician in Istanbul just uploaded. Instead of sending terabytes of hard drives across the country, the team uses Unity Catalog to securely share "Gold-level" datasets between government agencies and universities (like ITU or METU) for rapid disaster-risk assessment. 

\subsection{From Raw Waves to Real-Time Alerts: The Intelligent Data Tiering Through Medallion Architecture}
In geophysics, the Medallion Architecture on Databricks streamlines the transition from raw seismic sensor feeds to predictive insights by organizing data into three functional stages, Figure \ref{fig06}. The process begins with the Bronze layer, which ingests raw seismic waveforms from global APIs like the USGS or local station streams, preserving the original signal in its entirety. This data is then refined into the Silver layer, where signal processing removes "cultural noise" such as traffic or industrial vibrations to isolate tectonic activity. Finally, the Gold layer provides curated, analytics-ready datasets used for high-level earthquake magnitude and location prediction. This structured workflow allows research institutions to efficiently parallelize computationally intensive tasks like Template Matching, enabling the detection of micro-earthquakes hidden within years of continuous recording that traditional systems would miss.

\begin{figure}[H]
    \centering % Use \centering instead of the center environment to avoid extra vertical space
    \includegraphics[width=\textwidth]{Fig_06.png} 
    \caption{Geophysical Data Processing Framework: Istanbul EEW \& Seismic Analysis. Diagram source: authors.}
    \label{fig06}
\end{figure}

To foster a unified environment, the Databricks workspace is organized around a Medallion Architecture governed by Unity Catalog, ensuring a single source of truth across all tiers. By structuring data into Bronze (raw SEG-Y), Silver (physically validated), and Gold (model-ready features) catalogs, the organization eliminates manual data duplication and version chaos. Collaboration is further streamlined through Persona-Based Clusters—where engineers use memory-heavy nodes for ETL and scientists utilize GPU-enabled nodes for deep learning—and Databricks Repos, which integrates Git to ensure all code is versioned and peer-reviewed. Finally, replacing static PDF reports with Live SQL Dashboards allows geophysicists and scientists to interact directly with real-time seismic streams, effectively bridging the "Domain vs. Code" gap through a shared, transparent technical ecosystem.

By isolating the production environment, the mission-critical earthquake alert pipelines remain protected from experimental interference, while project-based folders in the development workspace unite engineers, geophysicists, and scientists within a single collaborative directory. This strategy, reinforced by a Medallion data architecture and persona-based compute clusters, eliminates version confusion and resource contention, allowing the team to transition seamlessly from raw seismic ingestion to real-time AI-driven insights within a unified, transparent ecosystem.

\subsection{Real-time seismic intelligence: A Medallion pipeline for the North Anatolian Fault}
A real-time monitoring system focused on the North Anatolian Fault (NAF) can include the workflow in Databricks with multi-layer structure.  Bronze Layer (Raw Ingestion) includes data collection. Here, the continuous 100Hz velocity streams from hundreds of stations (e.g., from AFAD’s network) ingested via APIs or direct satellite links. The data can be stored as raw binary or JSON blob format in Azure Data Lake Storage. Next is a physically validated Silver Layer which includes big data processing. At this stage, PySpark jobs applies Butterworth filters and removes noise from the Istanbul metro or heavy traffic on the E5 highway. As a result, the geophysicists obtain cleaned waveforms where the "P-wave" (primary) and "S-wave" (secondary) arrivals are sharp and visible. Finally, during the Gold Layer stage, the large datasets are undergone the predictive analytics and insights. This includes big data processing using ML. For instance, a DL model (CNN or LSTM) trained on historical Turkish earthquake catalogs (1975–2026) scans the Silver data. As a result, this enables real-time calculation of the earthquake’s epicenter, depth, and magnitude, which feeds directly into an early-warning dashboard for Istanbul.

\subsection{Governance and integration planning}

Governance and integration for this earthquake monitoring system must prioritize reliability, security, and interoperability by centralizing control through Unity Catalog to ensure a transparent, auditable lineage from sensor to alert. The implementation requires "hardened" security—utilizing Private Links and Identity Federation—to protect national seismic data, alongside Databricks Auto Loader for the resilient ingestion of high-frequency station streams. By integrating Spatial SQL for real-time risk mapping and using Delta Sharing to provide universities and government agencies with secure, instant access to "Gold" datasets, this framework transforms the Databricks environment into a high-trust, collaborative ecosystem capable of delivering rapid, scientifically-validated disaster response.

\section{Bridging Data Science and Geophysics: Collaborative Seismic Analytics for Enhanced Hazard Assessment}

\subsection{Transforming seismic data into real-time earthquake insights}

Databricks serves as a unified Data Intelligence Platform for geophysics by consolidating disparate seismic telemetry and sub-surface data into a single, collaborative environment for batch, streaming, and ML workloads. By centralizing high-throughput ingestion from global sensor networks alongside intensive seismic wave inversion and waveform processing, the platform eliminates the silos typically found between data architects and seismologists. This integration allows for real-time earthquake monitoring through automated pipelines that feed directly into predictive ML models for P-wave and S-wave analysis, all governed by a single security framework. Ultimately, this unified approach ensures a "single source of truth" for geophysical events, accelerating the transition from raw seismic noise to actionable early-warning hazard assessments with superior precision and scalability.

One of the important advantages in seismic big data processing through Databricks is that its notebooks facilitate seamless geophysical collaboration. This is ensured through multi-language support, enabling seismologists, data engineers, and analysts to execute Python, SQL, R, and Scala cells within a single environment. In a real-time monitoring workflow, this allows seismic engineers to ingest high-frequency waveform telemetry using Python, while seismologists simultaneously apply specialized R packages for statistical magnitude analysis or SQL to query historic earthquake catalogs. By using "magic commands" to switch languages, the team can analyze the same seismic data frames without exporting files, ensuring that complex waveform processing and hazard reporting remain synchronized in one unified, peer-reviewed research document.

\subsection{Enhancing Earthquake Early-Warning Systems Through Collective Computational Resources}

For geophysical monitoring teams, Databricks’ cluster sharing enables a unified computational theater where multiple researchers can concurrently process massive seismic telemetry streams using shared memory and resources. By attaching to a single high-performance cluster, data engineers and seismologists eliminate the "it works on my machine" problem, ensuring that complex waveform analysis and ML-driven hazard modeling are performed on a consistent, cost-effective infrastructure. This collective utilization of compute power facilitates rapid peer review and iterative model tuning, which is vital for the sub-second response times required in earthquake early-warning systems.

Databricks integrates with data lake storage to establish a unified seismic data lakehouse, utilizing direct connectivity to enable multi-disciplinary access to massive geophysical datasets without the overhead of data duplication. By leveraging data lake on top of the platform, it provides enables seismic data engineers to continuously stream high-frequency sensor telemetry into "Bronze" tables while seismologists concurrently perform non-destructive analysis and ML-driven phase picking on the same "Silver" and "Gold" datasets. Such architecture ensures a singular, version-controlled source of truth for waveform data, where schema enforcement and time-travel capabilities facilitate rigorous, reproducible research. Consequently, this integration eliminates the latency of manual data transfers, providing a scalable foundation for real-time earthquake monitoring and complex sub-surface modeling within a secure, collaborative framework.

\subsection{High-Performance Seismology with Apache and Azure Databricks Architecture}
The architecture of Apache and Azure Databricks provides solutions for big data handling in geophysics. Thus, it synthesizes high-performance seismic processing with enterprise-grade data governance, enabling teams to navigate the immense computational demands of modern seismology. By leveraging Apache Spark's distributed engine within the Azure cloud, researchers can execute massive parallel simulations, such as Full Waveform Inversion (FWI) or real-time earthquake arrival picking, across global sensor networks without the hardware constraints of traditional on-premise servers. This unified environment utilizes shared high-memory clusters and multi-language notebooks—where Python-based signal processing and SQL-driven cataloging coexist—to dissolve the boundaries between data architects and research seismologists. Furthermore, the seamless integration with Azure Data Lake Storage and Unity Catalog ensures that every terabyte of waveform telemetry is accessible through a secure, version-controlled "single source of truth," which is indispensable for reproducible scientific research and the sub-second latency required for critical early-warning hazard assessments.

\subsection{Implementation roadmap development}

The implementation roadmap follows a phased, value-driven approach that begins with a foundational pilot to establish a secure "Dev" environment and parallelize one high-value seismic network as a "shadow" pipeline. This is followed by an expansion phase, where historical seismic catalogs are migrated to enable large-scale AI training and geophysicists are onboarded through persona-based training. Finally, the system moves into operational maturity by deploying mission-critical early warning dashboards and automating CI/CD pipelines in the "Prod" environment, Figure \ref{fig07}. 

\begin{figure}[H]
    \begin{center}
	\hspace{-35pt}\includegraphics[width=13.0cm]{Fig_07.png}
    \end{center}
    \caption{Integrated CI/CD Pipeline for the Marmara Sea Early Warning System (EEW): From Seismic Code Repository to Live Production. Diagram source: authors.}
    \label{fig07}
\end{figure}

To minimize disruption, a "dual-run" strategy is used, maintaining legacy systems until the Databricks Gold layer is fully validated, while ensuring domain experts retain ownership of the scientific logic.

\subsection{Building a unified data platform for seismic stream integration in geophysics}

Optimizing a professional Databricks workspace for geophysics transforms fragmented seismic streams into a unified Data Intelligence Platform, essential for real-time earthquake monitoring and sub-surface modeling. By enforcing role-specific cluster policies—provisioning high-memory GPU nodes for intensive seismic wave inversion while restricting analysts to cost-effective SQL Warehouses—the environment balances computational rigor with fiscal governance. Granular Unity Catalog permissions ensure that data engineers maintain high-throughput ingestion pipelines from global sensor networks, while seismologists develop predictive ML models in isolated sandboxes using P-wave and S-wave arrival data. This architecture, fortified by Git-integrated notebooks for peer-reviewed code and automated CI/CD for waveform processing, eliminates silos between geophysicists and data architects. Ultimately, this structured ecosystem establishes a "single source of truth" for seismic events, accelerating the transition from raw telemetry to actionable hazard assessments with the precision and scalability required for critical early-warning systems.

%----------- section ----------->
\section{Deterministic Infrastructure and Performance Optimization for EEW Systems}

Programmatic cluster configuration and performance monitoring are essential for the technical reliability of Earthquake Early Warning (EEW) systems. Utilizing JSON-based configurations allows for the automated deployment of computational environments that are optimized for the low-latency requirements of seismic phase picking. This approach ensures that the production infrastructure is an exact, reproducible replica of validated research models, reducing the risk of manual configuration errors that could introduce processing delays during a seismic event. By defining resources like GPU-accelerated nodes programmatically, engineers ensure the system maintains the necessary throughput for high-frequency data analysis.

Complementing this, using Spark performance monitoring is necessary to maintain the stability of real-time telemetry streams. In an EEW context, monitoring for memory spills or CPU bottlenecks allows for the optimization of pipelines that must process sudden increases in data volume during earthquake swarms. Programmatic management of big data in geophysics supports the deployment of fault-tolerant, high-availability clusters capable of dynamic scaling as seismic activity propagates across a sensor network. Such technical framework ensures that the computational architecture remains stable and responsive, providing the consistent performance required for effective hazard detection and public safety.

\subsection{Computational Tiering and Latency Requirements in Earthquake Early Warning Workloads}

Earthquake Early Warning (EEW) systems rely on a tiered computational architecture characterized by increasingly intensive processing demands and stringent latency constraints. The primary workload involves the continuous, sub-second ingestion and filtering of high-velocity sensor telemetry of seismic signals, requiring deterministic I/O throughput to ensure raw data is available for immediate analysis. This is followed by event detection and phase picking, where machine learning models and signal processing algorithms must scale instantly to identify P-wave arrivals across thousands of concurrent streams during seismic swarms. 

While these real-time tasks prioritize minimal "data aging" to maximize warning lead times, they are complemented by compute-heavy, offline workloads such as Full Waveform Inversion (FWI), which utilize massive GPU clusters for complex geological modeling. Ultimately, the performance of an EEW system is defined by its ability to maintain high-availability streaming for life-safety detections while isolating these critical paths from the high-memory batch processes required for scientific refinement.

\subsection{Addressing Computational Limits in Seismic Early Warning Networks in Istanbul}

Effective Earthquake Early Warning (EEW) systems are primarily hindered by three critical technical bottlenecks that directly impact the "warning window" available to at-risk populations:

\begin{itemize}
\item data ingestion latency, 
\item computational saturation during seismic swarms, 
\item memory/shuffle bottlenecks of big data processing.
\end{itemize}

Data ingestion latency creates "aged" data, where every 100ms of delay in high-frequency telemetry streams reduces the geographic reach of life-saving alerts. Computational saturation during seismic swarms can paralyze a system if it lacks the ability to dynamically scale or utilize GPU acceleration for rapid phase picking, potentially causing the infrastructure to fail during peak demand. Finally, memory spills and shuffle bottlenecks within distributed computing frameworks like Spark can delay magnitude and location estimations, as shifting data between nodes during spatial-temporal correlations may force slow disk-based processing over rapid RAM-based calculations. Managing these issues through programmatic cluster configurations and rigorous performance monitoring is essential to maintain the sub-second speed and system credibility required for public safety.

\subsection{Hybrid Optimization Strategies for EEW in Istanbul}

In the high-stakes context of the Marmara Sea region, managing an EEW system for Istanbul requires a hybrid mosaic approach to balance extreme performance needs with budget realities. Maintaining the GPU-accelerated clusters necessary for sub-second phase picking is costly during seismic dormancy, yet essential for protecting critical infrastructure like the Marmaray tunnel and natural gas grids. To mitigate these pressures, engineers use programmatic configurations to offload "bursty" tasks—such as denoising earthquake swarms—to elastic cloud resources while keeping compute-intensive FWI on dedicated HPC clusters. Furthermore, implementing edge computing at nodal sensors across the Marmara Sea reduces data transfer bottlenecks and egress fees by transmitting only relevant signals. This integrated strategy, which includes monitoring for memory spills and CPU bottlenecks, ensures the computational architecture remains fiscally sustainable while delivering the deterministic speed required for public safety.

\subsection{GPU-Accelerated Systems for Marmara Region Earthquake Alerts}
To optimize seismic big data for the Istanbul and Marmara region, the strategy focuses on a tiered approach that balances high-speed performance with cost sustainability. By implementing edge-based signal processing at seafloor stations, the system filters raw noise at the source, reducing data transfer bottlenecks and egress fees. During seismic swarms, the architecture utilizes programmatic JSON configurations to trigger "warm" GPU-accelerated pools, preventing fatal memory spills and ensuring sub-second phase picking when it is most critical, Figure \ref{fig08}. 

\begin{figure}[H]
    \begin{center}
	\hspace{-35pt}\includegraphics[width=13.0cm]{Fig_08.png}
    \end{center}
    \caption{Earthquake Early Warning (EEW) system data flow on Databricks. Diagram source: authors.}
    \label{fig08}
\end{figure}

A hybrid mosaic model isolates high-priority alerting on dedicated HPC clusters while offloading bursty, offline tasks like Full Waveform Inversion (FWI) to elastic cloud resources. Finally, the integration of Physics-Informed Neural Networks (PINNs) allows for greater geological accuracy with fewer compute cycles, ensuring the system can provide up to 20 seconds of warning for the Anatolian side of Istanbul while remaining fiscally viable.

\subsection{Hybrid Cluster Architecture for Sub-Second Earthquake Detection and Alerts}

In the Marmara Sea region, an EEW-optimized Databricks cluster functions as a mission-critical engine within a hybrid mosaic architecture. Azure Databricks serves as a managed, high-performance platform built directly upon the open-source Apache Spark engine. In the seismic monitoring of Istanbul region, Spark provides the core distributed computing power necessary to parallelize seafloor telemetry data across a cluster, while Databricks enhances this with a managed runtime that optimizes memory management and I/O throughput. Such integration transforms Spark's raw computational capabilities into a collaborative, mission-critical environment for Istanbul’s defense, automating the complex orchestration of cluster resources and providing a unified workspace for developing real-time earthquake early warning (EEW) algorithms.

By utilizing programmatic JSON configurations and GPU-accelerated pools, the system handles "seismic swarms" with sub-second latency, preventing fatal memory spills during the "Golden Seconds". The workflow integrates seafloor "smart nodes" for edge-based ingestion (Bronze), Physics-Informed Neural Networks for real-time denoising (Silver), and refined "Digital Twin" models for hazard mitigation (Gold). This systematic optimization balances life-saving performance with fiscal sustainability, offloading compute-intensive inversions to elastic cloud resources while maintaining dedicated high-availability clusters for Istanbul’s public safety.

\subsection{Centralized Control of Hybrid Clusters in EEW Systems in Istanbul}

In the seismotectonic framework of the Marmara region, the master node (or Driver) serves as the primary orchestrator within a distributed, high-concurrency computing architecture. Operating as the "central nervous system" of a hybrid mosaic cluster, its role is to decompose complex geophysical algorithms—such as automated phase picking and magnitude estimation—into discrete, parallelizable tasks assigned to a network of worker nodes. By ingesting high-frequency telemetry from seafloor "smart nodes" and terrestrial stations, the master node ensures sub-second latency through optimized Directed Acyclic Graph (DAG) scheduling. 

In the event of a high-magnitude seismic swarm, the master node maintains the integrity of the EEW pipeline by dynamically redistributing workloads and managing cluster resources in real-time. This structural oversight facilitates the rapid transformation of raw waveform data into deterministic hazard alerts, ensuring that the "Golden Seconds" of lead time are preserved for Istanbul’s critical infrastructure and public safety protocols. hence, the distributed computing's primary advantage in the Marmara region is its ability to process seismic datasets that far exceed the memory capacity of any individual machine, ensuring the computational "headroom" necessary for Istanbul’s safety. 

By leveraging horizontal scalability, the architecture partitions massive telemetry streams of large massifs of seismic data across a hybrid mosaic of nodes, enabling the parallel processing required to achieve sub-second phase picking during the "Golden Seconds" of an EEW. This distributed approach eliminates the risk of fatal memory spills to disk, allowing complex tasks like FWI to run efficiently across a cluster rather than hitting the physical ceiling of a single server. Ultimately, this parallel power transforms the high-velocity Big Data from seafloor sensors into actionable alerts, maintaining the deterministic speed and reliability required to protect critical infrastructure.

\section{High-Performance Seismic Data Storage for EEW in Marmara Region}

\subsection{Strategic Alignment Analysis for Seismic Big Data in Business Context}

For geophysical organizations operating in the Marmara region in Istanbul with centralized control in Ankara, strategic cloud alignment must prioritize ultra-low latency and seismic resilience to maximize the narrow warning window for Istanbul's high-density population. The primary business objective is establishing a deterministic edge-to-cloud architecture that can process high-frequency backhaul from the North Anatolian Fault in milliseconds, ensuring EEW alerts reach the metropolis before the destructive S-waves. This requires a provider with geographically redundant regions located outside the primary fault zone to maintain operational integrity during a "Big One," combined with an "open science" data framework to facilitate rapid sharing of MiniSEED data with international partners. By leveraging spot instance orchestration for large-scale ambient noise tomography and high-throughput object storage for the region's massive historical catalogs, the geophysical data centres in Istanbul can balance the high-performance demands of real-time public safety with the cost-efficiency required for long-term seismological research.

Finding a single cloud provider to anchor an Earthquake Early Warning (EEW) system for the Marmara region involves a fundamental tension between centralized cloud efficiency and localized seismic physics. The challenges are primarily driven by the proximity of infrastructure to the North Anatolian Fault and the rigid latency requirements. The necessity for sub-second processing and transmission is defined by the stringent temporal constraints of the earthquake early warning (EEW) latency window, frequently referred to in seismic risk mitigation as the "Golden Seconds." This critical interval represents the brief period between the detection of initial P-waves and the arrival of destructive S-waves, requiring deterministic, near-real-time computational throughput to facilitate effective automated response and public notification.

\subsection{Decoupling Storage and Compute for Future Geophysical Data Challenges in EEW}

To maintain long-term strategic positioning in EEW within the Marmara region, a geophysical organization must prioritize an evolutionary architecture of big data storage that decouples storage from compute, specifically to accommodate the massive "step-function" data growth generated by Distributed Acoustic Sensing (DAS) along the North Anatolian Fault. By standardizing on open-source data formats (e.g., Zarr or Parquet) and adopting silicon-agnostic containerization, the organization ensures that critical earthquake early warning (EEW) algorithms can seamlessly migrate to future, high-efficiency hardware such as TPUs or FPGAs. Furthermore, utilizing hybrid-cloud management frameworks ensures that Istanbul’s monitoring infrastructure remains both technologically cutting-edge and operationally resilient, allowing for the localized, low-latency processing required for real-time public safety while leveraging the cloud for data-intensive, AI-driven seismic research.

\section{Decision Framework Development in Geophysical Big Data}

\subsection{A Tiered Framework for Life-Safety in Marmara Seismic Cloud Systems}

The assessment of cloud infrastructure for seismic monitoring in the Marmara region necessitates a multidimensional tiered framework that prioritizes life-safety strategy in seismically active region of Türkiye. This approach utilizes a technical determinism gate, requiring a p99.9 latency of less than 50ms for P-wave detection and the geographic isolation of failover regions from the North Anatolian Fault to ensure seismic resilience. 

Big data infrastructure viability for seismic data storage is evaluated through an operational synergy tier. The focus is on integration of big data and AI for national security auditing compliance, alongside a strategic agility tier designed to accommodate "step-function" data expansion from Distributed Acoustic Sensing (DAS) via open-source standards and containerization. Finally, an economic total value analysis integrates cost-optimization strategies—including reserved capacity for baseline monitoring and high-performance computations—to balance the exigencies of the "emergent seismic warning with long-term data storage sustainability.

\subsection{Combining Edge Filtering and Serverless Computing for Seismic Analysis}

By automating lifecycle policies to migrate historical waveforms to deep-archive tiers (e.g., Google Cloud Archive Storage, Azure Archive or AWS Glacier Deep Archive), seismic data centres can realize a significant reduction in storage fees while maintaining long-term research data for fault-line analysis in tectonically active zones of Marmara Sea. Simultaneously, transitioning from "always-on" servers to Spot Instances or serverless triggers for seismic event detection cuts compute costs. When combined with Edge Filtering to suppress noise at the sensor site, these techniques drastically reduce hidden egress fees, allowing the organization to pivot its budget from infrastructure maintenance to critical real-time earthquake early warning innovation.

For geophysical organizations in Istanbul, the most significant opportunities for big data workflow optimization lie in transitioning from monolithic, "always-on" infrastructure to a highly elastic, event-driven architecture. By implementing a high computing strategy—leveraging Spot instances or serverless engines to process Terabytes of daily data in high-intensity bursts rather than 24/7 cycles—seismic data centres can reduce compute expenses. Furthermore, priorities in big data processing are in automated storage tiering, moving massive seismic archives to deep-cold tiers, and deploying edge computing to filter background noise at the sensor site, which drastically cuts data egress and ingress fees. Finally, consolidating fragmented security and compliance tools into a unified governance frameworks eliminates redundant Software as a Service (SaaS) licensing and the costs of engineering labor, allowing the organization to pivot budget away from maintenance toward real-time early warning innovation.

\subsection{Operational Constraints for Istanbul’s Seismic Monitoring Infrastructures}
The optimization of Istanbul’s geophysical data architectures is strictly delimited by the latency-critical imperatives of the Marmara Sea Early Warning System (EEW). Because the North Anatolian Fault’s proximity demands sub-second P-wave processing, the use of high-latency archival tiers is technically precluded for active telemetry, requiring high-cost "hot" storage to ensure alert delivery before S-wave arrival. Finally, the requirement for system durability during peak tectonic events restricts the deployment of volatile Spot instances, mandating a baseline of redundant, high-availability compute. Consequently, the cloud strategy for the Istanbul metropolitan region must prioritize operational resilience and public safety over raw fiscal reduction, transforming cost management into a complex multi-objective balancing task.

\subsection{Balancing Latency and Durability in Istanbul’s Seismic Data Systems Through Phased Implementation Framework}

The optimization of Istanbul’s geophysical infrastructure through seismically resilient and cost-optimized geophysical architectures for big data storage follows a tiered implementation timeline that prioritizes high-yield, low-risk fiscal recovery before transitioning to complex architectural hardening. 

\begin{itemize}
\item Phase 1 focuses on Intelligent Storage Tiering, leveraging automated lifecycle policies to migrate petabytes of legacy seismic archives to deep-cold tiers, yielding an immediate reduction in storage overhead without impacting real-time telemetry. 
\item Phase 2 introduces Hybrid Compute Orchestration, balancing a baseline of Reserved Instances for mission-critical P-wave detection with elastic Spot instances for batch-oriented research, effectively slashing compute costs while maintaining the durability required for Marmara Sea resilience.  
\item Phase 3 implements Edge Processing at sensor sites along the North Anatolian Fault to filter seismic noise locally, significantly suppressing network egress fees and latency. 
\end{itemize}

This strategic workflow ensures that initial administrative savings self-fund the subsequent hardware and DevOps innovations necessary to satisfy both KVKK (Kişisel Verileri Korunması Kanunu) data sovereignty and the life-safety mandates of the Istanbul metropolitan region.

\subsection{Risk-Weighted Optimization Strategy for EEW in the Marmara Sea}
Balancing cost reduction with the life-safety mandates of the Marmara Sea EEW requires a risk-weighted optimization strategy that treats sub-second latency as a non-negotiable "performance floor." This is achieved by system decoupling, where the mission-critical P-wave detection pipeline is isolated on high-performance, reserved infrastructure to ensure 99.99\% durability, while non-essential research and archival workloads are aggressively offloaded to volatile Spot instances and deep-cold data storage. 

To further reduce costs without sacrificing performance, Edge-heavy decentralization is employed to perform initial seismic "picking" at the sensor site, suppressing raw data volumes by over 90\% and reducing expensive cloud egress and ingestion fees. Finally, the implementation of trigger-based burst scaling ensures that the organization maintains a low-cost "Pilot Light" footprint during seismic quiescence, instantly scaling to high-capacity compute only when a confirmed detection event occurs. This multi-tiered approach maintains KVKK compliance and public safety standards while systematically eliminating the "idle capacity" that typically inflates geophysical operational budgets.

For instance, the six-month implementation roadmap for Istanbul’s geophysical infrastructure is a phased strategy that balances immediate fiscal recovery with long-term seismic resilience. Months 1–2 prioritize financial stabilization through the migration of legacy seismic archives to deep-cold storage and the alignment of data inventories with KVKK residency mandates, yielding immediate OpEx reduction. Months 3-4 transition to architectural hardening, establishing a "Pilot Light" model that utilizes Reserved Instances for 24/7 P-wave detection while offloading non-critical tomographic modeling to volatile Spot fleets. The timeline concludes in Months 5–6 with the deployment of Edge Processing units at Marmara Sea sensor stations to filter seismic noise locally, coupled with in-region cloud consolidation to eliminate international egress fees. This structured approach ensures that initial storage savings fund the hardware and DevOps innovations required to minimize the early-warning blind zone and maintain 99.99\% system durability during peak tectonic events.

\subsection{AI-Driven Framework for Managing Cloud Costs in Seismic Monitoring}
To mitigate the risk of "cloud sprawl" following the initial deployment, the framework employs AI-driven anomaly detection to identify fiscal deviations during high-frequency seismic swarms, triggering automated right-sizing reviews. This technical governance is reinforced by a structured operational cadence, prioritizing Reserved Instance (RI) optimization to ensure that ca. 90\% of the baseline detection workload is sustained on high-availability, discounted infrastructure. By embedding cost-impact simulations directly into the geophysical Continuous Integration and Continuous Deployment CI/CD pipeline, the model transforms regulatory and fiscal oversight into an automated, invisible standard, thereby ensuring that operational savings are perpetually redirected toward enhancing the spatiotemporal accuracy of the North Anatolian Fault monitoring network, Figure \ref{fig05}.

The successful execution of change management and stakeholder alignment for the Marmara Sea EEW necessitates a transition from traditional CapEx-heavy procurement to a decentralized, FinOps-oriented OpEx model, reconciling the divergence between research-driven agility and mission-critical operational stability. Organizational integration is achieved by implementing "Golden Path" self-service architectures that empower geophysicists while programmatically enforcing KVKK data sovereignty and fiscal guardrails through PaC. By mitigating "automation anxiety" through incremental CI/CD validation—benchmarked against historical NAF seismotectonics—the framework fosters a collaborative ecosystem where legal, financial, and technical stakeholders utilize automated monitoring loops to maintain system resilience and spatiotemporal detection accuracy.

A robust cost optimization governance framework for the Marmara Sea EEW must utilize a proactive, automated, and scalable architecture to preserve the fiscal integrity of high-frequency geophysical data streams. This is achieved through the integration of real-time cost anomaly detection—acting as a financial seismograph to identify "runaway" computational processes—and automated threshold responses that programmatically deprovision non-essential development environments while maintaining the sanctity of the mission-critical seismic monitoring path. By instituting a tiered notification system with progressive budget alerts and formal escalation procedures, the framework ensures that technical and administrative stakeholders possess sufficient lead time for strategic resource reallocation, effectively mitigating the risk of budgetary tremors without compromising the sub-second latency required for regional public safety.

\subsection{Elastic Scaling and PaC in Istanbul’s EEW Framework}
The evaluation of geophysical architectural optimization for the Istanbul megacity must transcend rudimentary fiscal metrics, prioritizing seismological resilience and governance maturity within the context of the high-risk Marmara seismic gap. Technical efficacy is benchmarked against the rigorous spatiotemporal constraints of the North Anatolian Fault (NAF), where success is defined by maintaining a "detection-to-dissemination" p99 latency threshold of $\leq 200$ ms—a critical requirement for mitigating the "Blind Zone" during a potential $Mw \geq 7.0$ event. Following the catastrophic 2023 Kahramanmaraş earthquakes, which underscored the necessity for robust, non-volatile telemetry under extreme tectonic stress, this framework mandates 100\% system durability during high-frequency seismic swarms through FinOps-driven elastic scaling.

Furthermore, the big data governance strategy enforces absolute KVKK data sovereignty, utilizing Policy-as-Code (PaC) "Golden Paths" to prevent trans-border data egress of sensitive telemetry while ensuring audit-readiness for national critical infrastructure. Ultimately, success is quantified through the Unit Cost of Detection, an efficiency ratio that aims to minimize the cost-per-event through edge-processing and noise-filtering, thereby optimizing the fiscal-to-safety conversion rate for the Marmara Sea EEW system.

The sustained efficacy of the proposed cloud optimization framework for the Marmara Sea Early Warning System (EEW) is predicated upon a closed-loop governance architecture that integrates Automated Guardrails, FinOps telemetry, and PaC. By implementing PaC protocols, the infrastructure programmatically enforces KVKK-compliant data residency, utilizing deterministic region-locking to prevent unauthorized trans-border data egress. This architectural rigidity is balanced by continuous visibility via specialized dashboards that monitor the Unit Cost of Detection—a critical performance indicator that normalizes computational expenditure against seismic event frequency.

%----------- section ----------->
\section{Conclusions}

When we think of big data in seismology, the first thing that comes to mind is not a collection of spreadsheets or databases—it is patterns. Data obtained from large massifs of seismic signals record a digital pattern of geological processing happening in a particular place in the world and not visible to human eyes. On its own, a single data point in seismic record is barely a lonely number in a table. But when we have the large collection of those records in large quantities, those points connect to tell a story about the earthquakes, the behavior of inner mantle, the geophysical processes, or regional trends in earthquake activities. The challenge of big data comes with the transform of those records to big data due to their volume. For seismic domain, we are now in the era of Petabytes (1,000 Terabytes) and even Exabytes (1,000 Petabytes) of record sources. Consequently, novel methods of big data processing are emerging. 

Nowadays, the successful implementation of cloud optimization for the Istanbul EEW requires a multi-layered communication strategy that reframes technical and fiscal guardrails as essential enablers of seismic resilience. By categorizing benefits into FinOps-driven fiscal predictability for leadership, "Golden Path" technical agility for researchers, and automated Policy-as-Code (PaC) compliance for legal stakeholders, the organization can resolve the inherent tension between rapid innovation and strict KVKK data sovereignty. Potential anxieties regarding automated cost controls are mitigated by establishing a "Mission-Critical Whitelist" that prioritizes life-safety alerts over cost reduction, utilizing "Shadow Mode" to demonstrate reliability before activation, and maintaining a CI/CD pipeline that transitions from human-in-the-loop delivery to fully automated deployment only after rigorous validation against historical North Anatolian Fault data. Ultimately, this approach creates a self-funding infrastructure cycle where operational savings are systematically reinvested into the high-performance computing assets required for the next generation of earthquake forecasting.

In a geophysical context of EEW in Marmara, success of big data processing is not merely defined by the volume of petabytes archived, but by the system's ability to maintain a high Unit Cost of Detection. Hence, the optimization must significantly lower long-term storage overhead without introducing prohibitive latency that would impede the ability of the scientific community of geophysicists to analyze the North Anatolian Fault’s behavior following a major seismic event. The primary metric for evaluating the efficacy of big data lifecycle strategy in Istanbul region is the maximization of fiscal efficiency. This can be achieved by transitioning high-volume seismic waveforms to lower-cost storage tiers, calibrated against the mission-critical requirement for rapid data retrieval during post-event analysis or seismotectonic model validation of the Marmara Sea region. 


%----------- section ----------->
\authorcontributions{Supervision, conceptualization, administration, funding acquisition, and project administration -- C. A. Z. Methodology, software, resources, funding acquisition, writing---original draft preparation, methodology, data curation, visualization, formal analysis, validation, writing---review and editing, and investigation, P.L.}

%\funding{The publication was funded by the Editorial Office of XXXX, Multidisciplinary Digital Publishing Institute (MDPI), by providing 100\% discount for the APC of this manuscript.}

\institutionalreview{Not applicable.}

\informedconsent{Not applicable.}

\dataavailability{Not applicable} 

\acknowledgments{The authors thank the reviewers for reading and review of this manuscript.}

\conflictsofinterest{The authors declare no conflict of interest.} 

\abbreviations{Abbreviations}{
The following abbreviations are used in this manuscript:\\

\noindent 
\begin{tabular}{@{}ll}
AFAD & Disaster and Emergency Management Authority \\
AI & Artificial Inteligence \\
API & Application Programming Interface \\
CapEx & Capital Expenditure \\
CCS & Carbon Capture and Storage \\
CNN & Convolutional Neural Networks\\
CSV & Comma Separated Value\\
DAS & Distributed Acoustic Sensing \\
DAG & Directed Acyclic Graph \\
DL & Deep Learning \\
EEW & Earthquake Early Warning \\
ETL & Extract, Transform, and Load \\
FWI & Full Waveform Inversion \\
GPU & Graphics Processing Unit \\
GNNs & Graph Neural Networks \\
HPC & High-Performance Computing \\
I/O & Input-Output \\
KOERI & Kandilli Observatory and Earthquake Research Institute\\
KVKK & Kişisel Verileri Korunması Kanunu \\
LSTM & Long Short-Term Memory \\
ML & Machine Learning \\
NAF & North Anatolian Fault \\
NEMC & National Earthquake Monitoring Center \\
OpEx & Operating Expenditure \\
PINNs & Physics-Informed Neural Networks \\
PDF & Portable Document Format \\
PaC & Policy-as-Code \\
QC & Quality Control \\
RNNs & Recursive Neural Networks \\
SaaS & Software as a Service \\
SEG-Y & Society of Exploration Geophysicists - Y format \\
SHRA & Seismic Hazard Risk Assessment \\
SQL & Structured Query Language \\
SNR & Signal-to-Noise Ratio \\
\end{tabular}
}

%%%%%%%%%%%%%%%%%%%%%%%%%%%%%%%%%%%%%%%%%%
\begin{adjustwidth}{-\extralength}{0cm}
%\printendnotes[custom] % Un-comment to print a list of endnotes

\reftitle{References}
\nocite{*}

%=====================================
% References, variant A: external bibliography
%=====================================
\bibliography{Bibliography.bib}

%%%%%%%%%%%%%%%%%%%%%%%%%%%%%%%%%%%%%%%%%%
\end{adjustwidth}
\end{document}
